%
% 2-liouville.tex
%
% (c) 2026 Prof Dr Andreas Müller
%
\section{Das Prinzip von Liouville
\label{buch:intalg:section:liouville}}
\kopfrechts{Das Prinzip von Liouville}%
Die in Kapitel~\ref{chapter:ratint} durchgerechneten Beispiele 
uggerieren, dass eine Stammfunktion immer nur eine Linearkombination
von Logarithmen der Teilnenner der Partialbruchzerlegung ist.
Das Prinzip von Liouville zeigt, dass dies ganz allgemein gilt.

%
% Das Prinzip
%
\subsection{Das Prinzip 
\label{buch:intalg:liouville:subsection:prinzip}}

\begin{satz}
Sei $F$ ein Funktionenkörper mit Konstantenkörper $K$ und $f\in F$.
Wenn $f$ eine Stammfunktion $g$ in einer elementaren Erweiterung
$G\supset F$ hat, die den gleichen Konstantenkörper wie $F$ hat, dann
gibt es $v_0,v_1,\dots,v_n\in F$ und $c_1,\dots,c_k\in K$ mit
\begin{equation}
f
=
Dv_0
+
\sum_{i=1}^n c_i\frac{Dv_i}{v_i},
\end{equation}
d.~h.~
\[
\int f(z)\,dz
=
v_0
+ 
\sum_{i=1}^n c_i \log(v_i).
\]
$G$ muss also alle Logarithmen $\log(v_i)$, $i=1,\dots,n$, enthalten.
\end{satz}

Wie kann man überhaupt eine solche Aussage beweisen?
Wenn $g$ die Stammfunktion von $f$ ist, dann muss die Ableitung
von $g$ wieder $f$ ergeben, also $Dg=f$.
Wenn also $g$ in einer Erweiterung $F(\vartheta_1,\dots,\vartheta_n)$
von $F$ liegt, dann kann man versuchen, einen Ausdruck für $g$ in den
Elementen $\vartheta_k$ abzuleiten.
Dabei muss ein Ausdruck entstehen, welcher mit $f$ übereinstimmt.
Alle Terme mit den Elementen $\vartheta_k$ müssen sich daher wegheben.

Der allgemeine Fall einer Erweiterung der Form
$F(\vartheta_1,\dots,\vartheta_n)$ führt zu ziemlich komplizierten
Ausdrücken.
Daher beginnen wir damit einer Erweiterung $F(\vartheta)$ um ein
einziges $\vartheta$.
In diesem Fall muss $g$ die Form einer rationalen Funktion
\begin{equation}
g
=
\frac{a(\vartheta)}{b(\vartheta)} 
=
\frac{
a_0 + a_1\vartheta + \dots a_n\vartheta^n
}{
b_0 + b_1\vartheta + \dots b_m\vartheta^m
},
\label{buch:intalg:liouville:eqn:bruch}
\end{equation}
wobei $a(\vartheta)\in F[\vartheta]$ und $b(\vartheta)\in F[\vartheta]$
Polynome mit Koeffizienten in $F$ sind.
In diesem Fall lässt sich die Ableitung von $g$ leichter 
berechnen.
Zu diesem Zweck berechnen wir im
Abschnitt~\ref{buch:intalt:liouville:subsection:polynome}
zunächst die Ableitung von Polynomen von $\vartheta$ für die
verschiedenen Arten von Erweiterungen.
Diese Resultate verwenden wir dann im
Abschnitt~\ref{buch:intalt:liouville:subsection:spezialfaelle},
um das Prinzip von Liouville für eine Erweiterung der Form $F(\vartheta)$
zu beweisen.
In Abschnitt~\ref{buch:intalt:liouville:subsection:allgemein}
wird dann der allgemeine Fall bewiesen.

%
% Polynome von Erweiterungen
%
\subsection{Polynome von Erweiterungen
\label{buch:intalg:liouville:subsection:polynome}}
Die oben formulierte Beweisstrategie für das Prinzip von Liouville
basiert darauf, dass wir die Ableitung der Stammfunktion ausrechnen
und dann nachweisen, dass nur logarithmische Terme als Ableitung
Ausdrücke in $F$ ergeben.
Dazu müssen zunächst die Ableitungen von Zähler und Nenner in
\eqref{buch:intalg:liouville:eqn:bruch}
verstanden werden, was in den folgenden Abschnitten geklärt werden
soll.

%
% Polynome einer logarithmischen Erweiterung
%
\subsubsection{Polynome einer logarithmischen Erweiterung}
Das Prinzip von Liouville verspricht, dass nur logarithmische Erweiterungen
für die Stammfunktion gebraucht werden, und dass diese nur linear
vorkommen.
Dies wird sich darin äussern, dass die Erweiterung im rationalen Teil
nicht vorkommen kann.
Dies wird sich daraus ergeben, wie genau sich der Grad eines Polynoms
in $\vartheta$ beim Ableiten ändert.
Der folgende Satz gibt darüber Auskunft.

\begin{satz}
Sei $F$ ein Funktionenkörper mit Konstantenkörper $K\subset F$ und sei
$\vartheta$ ein logarithmisch transzendentes Element über $F$ mit
$\vartheta=\log u$ mit $u\in F$.
Sei ausserdem
$a(\vartheta) = a_0+a_1\vartheta+\dots+a_n\vartheta^n\in F[\vartheta]$
ein Polynom in $\vartheta$ mit Koeffizienten in $F$.
Dann gilt:
\begin{enumerate}
\item
Die Ableitung von $a(\vartheta)$ ist ebenfalls ein Polynom mit Koeffizienten
in $F$: $Da(\vartheta)\in F[\vartheta]$.
\item
Wenn der Leitkoeffizient $a_n$ eine Konstante ist, dann ist
$\deg Da(\vartheta)=\deg a(\vartheta) -1$.
\item
Wenn der Leitkoeffizient $a_n$ keine Konstante ist, dann ist
$\deg Da(\vartheta)=\deg a(\vartheta)$.
\end{enumerate}
\end{satz}

\begin{proof}
\begin{enumerate}
\item
Die Ableitung ist
\begin{align*}
Da(\vartheta)
&=
D\sum_{k=0}^n a_k\vartheta^k
=
\sum_{k=0}^n \bigl(Da_k\vartheta^k + a_kk\vartheta^{k-1}D\vartheta\bigr)
\\
&=
\sum_{k=0}^n Da_k\vartheta^k
+
D\vartheta \sum_{k=0}^{n-1} a_{k+1}(k+1)\vartheta^k
\\
&=
\sum_{k=0}^{n-1} \bigl(Da_k+(k+1)a_{k+1}D\vartheta\bigr) \vartheta^k
+
Da_n \vartheta^n
\\
&=
\sum_{k=0}^{n-1}
\biggl(
\underbrace{Da_k+(k+1)a_{k+1}\frac{Du}{u}}_{\displaystyle\in F}
\biggr) \vartheta^k
+
Da_n \vartheta^n
\end{align*}
Da $Du\in F$ und $u\in F$ folgt, dass der Koeffizient in der Klammer
ebenfalls in $F$ ist und daher ist $Da(\vartheta)\in F[\vartheta]$.
\item
Wenn $a_n$ eine Konstante ist, dann fällt der Term mit $\vartheta^n$ weg,
der Grad von $Da(\vartheta)$ ist also höchstens $n-1$.
Es ist zu zeigen, dass der Grad genau $n-1$ ist, d.~h.~dass der Koeffizient
von $\vartheta^{n-1}$ nicht verschwindet.
Würde dieser Koeffizient verschwinden, dann wäre wegen $Da_n=0$
\begin{align*}
0
&=
Da_{n-1} + na_n D\vartheta
\\
&=
D(a_{n-1}+na_n\vartheta),
\end{align*}
d.~h.~$a_{n-1}+na_n\vartheta=c$ ist eine Konstante $c\in K$.
Es folgt die Gleichung
\[
na_n\cdot\vartheta
+
(\underbrace{a_{n-1}-c}_{\displaystyle\in K})
=
0
\]
bedeutet aber, dass $\vartheta$ algebraisch ist über $K$.
Dies widerspricht der Voraussetzung, dass $\vartheta$ transzendent ist.
Der Widerspruch zeigt, dass $\deg Da(\vartheta)=n-1$ ist.
\item
Wenn $a_n$ keine Konstante ist, dann ist $Da_n\ne 0$ und damit
$\deg Da(\vartheta)=n$.
\qedhere
\end{enumerate}
\end{proof}

%
% Polynome einer exponentiellen Erweiterung
%
\subsubsection{Polynome einer exponentiellen Erweiterung}
Im Gegensatz zu einer logarithmischen Erweiterung nimmt der Grad
eines Polynoms einer exponentiellen Erweiterung beim Ableiten nicht ab.

\begin{satz}
Sei $F$ ein Funktionenkörper mit Konstantenkörper $K\subset F$ und sei
$\vartheta$ ein exponentielles transzendentes Element über $F$ mit
$\vartheta=\exp u$ mit $u\in F$.
Sei ausserdem
$a(\vartheta) = a_0+a_1\vartheta+\dots+a_n\vartheta^n\in F[\vartheta]$
ein Polynom in $\vartheta$ mit Koeffizienten in $F$.
Dann gilt:
\begin{enumerate}
\item
Für $h\in F$ und $n\ne 0$ gilt $D(h\vartheta^n) = \bar{h}\vartheta^n$ mit
$\bar{h}\in F$ und $\bar{h}\ne 0$.
\item
Für ein Polynom
$a(\vartheta)\in F[\vartheta]$
gilt
$\deg Da(\vartheta) = \deg a(\vartheta)$.
\item
$a(\vartheta) \in F[\vartheta]$
teilt $Da(\vartheta)$ genau dann, wenn $a(\vartheta)$ ein Monom ist,
also $a(\vartheta)=h\vartheta^n$ mit $h\in F\setminus\{0\}$, $n\in\mathbb{Z}$.
\end{enumerate}
\end{satz}

\begin{proof}
\begin{enumerate}
\item
Die Ableitung ist
\[
D(h\vartheta^n)
=
Dh\cdot\vartheta^n + hn\vartheta^{n-1}D\vartheta
=
Dh\cdot\vartheta^n + hn\vartheta^{n-1} \vartheta Du
=
(\underbrace{Dh+ hn Du}_{\displaystyle=\bar{h}}) \vartheta^n
\]
Wäre $\bar{h}=0$, dann würde gelten $h\vartheta^n$ eine Konstante $c\in K$,
d.~h.~es gilt eine Gleichung der Form
\[
h\cdot \vartheta^n - c = 0.
\]
Dies bedeutet, dass $\vartheta$ algebraisch ist über $F$, im Widerspruch
zur Voraussetzung, dass $\vartheta$ transzendent ist.
Somit folgt $\bar{h}\ne 0$.
\item
Die Ableitung des Terms $a_n\vartheta^n$ des Polynoms ist
$\bar{a}_n\vartheta^n$ mit $\bar{a}_n\ne 0$.
Somit ist $\deg Da(\vartheta) = n = \deg a(\vartheta)$.
\item
Wenn $a(\vartheta)$ das Monom $h\vartheta^n$ ist, dann ist
\[
Da(\vartheta)
=
\bar{h}\vartheta^n
=
\frac{\bar{h}}{h}
\cdot
h\vartheta^n
=
\frac{\bar{h}}{h}
\cdot
a(\vartheta),
\]
d.~h.~wegen $\bar{h}/h\in F$ ist $Da(\vartheta)$ ist ein Teiler von
$a(\vartheta)$.

Sei jetzt umgekehrt $Da(\vartheta)$ ein Teiler von $a(\vartheta)$.
Es gibt dann ein $g\in F$ mit $Da(\vartheta)=ga(\vartheta)$.
Sei $m$ der höchste Exponent $<n$ von $\vartheta$ im Polynom $a(\vartheta)$,
welcher einen nicht verschwindenen Koeffizienten $a_m\ne 0$ hat.
Das Polynom hat also die Form
\begin{align*}
a(\vartheta)
&=
a_n\vartheta^n + a_m\vartheta^m + \ldots + a_1\vartheta + a_0
\intertext{mit der Ableitung}
Da(\vartheta)
&=
\bar{a}_n\vartheta^n
+
\bar{a}_m\vartheta^m
+
\ldots
+
\bar{a}_1\vartheta
+
Da_0.
\end{align*}
mit $\bar{a}_k=Da_k+ka_kDu$.
Die Teilbarkeitsbedingung bedeutet dann
\[
\renewcommand{\arraycolsep}{2pt}
Da(\vartheta)
=
ga(\vartheta)
\quad
\Rightarrow
\quad
\left\{\;
\begin{array}{lclclcl}
ga_n &=& \bar{a}_n &=& Da_n &+& na_n Du \\
ga_m &=& \bar{a}_m &=& Da_m &+& ma_m Du.
\end{array}
\right.
\]
Auflösen nach $g$ ergibt
\begin{align*}
g
&=
\frac{Da_n}{a_n} + n\,Du
=
\frac{Da_m}{a_m} + m\,Du.
\intertext{Die Differenz der rechten Seiten ist}
0
&=
\frac{Da_n}{a_n}
-
\frac{Da_m}{a_m}
+
(n-m)\,Du
\intertext{Die rechte Seite kann man auch als}
0
&=
D
\biggl(
\frac{a_n}{a_m}
+
(n-m)u
\biggr)
\end{align*}
Die rechte Seite ermöglicht auch, die Ableitung
\begin{align*}
D
\biggl(
\frac{a_n}{a_m}\vartheta^{n-m}
\biggr)
&=
\frac{(Da_n) a_m - a_n\,Da_m}{a_m^2} \vartheta^{n-m}
+
(n-m)\vartheta^{n-m}Du
\\
&=
\frac{a_n}{a_m}
\biggl(
\frac{Da_n}{a_n}
-
\frac{Da_m}{a_m}
-
(n-m)Du
\biggr)
=
0
\end{align*}
zu berechnen.
Somit ist $(a_n/a_m)\vartheta^{n-m}$ eine Konstante, was nur für $n-m=0$
möglich ist.
Es gibt also gar keinen Koeffizienten $a_m\ne 0$ mit $m<n$.
Somit ist $a(\vartheta)=a_n\vartheta^n$ ein Monom.
\qedhere
\end{enumerate}
\end{proof}

%
% Polynome einer algebraischen Erweiterung
%
\subsubsection{Polynome einer algebraischen Erweiterung}

%
% Beweis des Prinzips von Liouville in Spezialfällen
%
\subsection{Beweis des Prinzips von Liouville in Spezialfällen
\label{buch:intalg:liouville:subsection:spezialfaelle}}
Das Prinzip von Liouville geht davon aus, dass eine Funktion $f$
in einem Funktionenkörper eine Stammfunktion $g$ in einer geeigneten
Erweiterung hat.
Es zeigt dann, dass diese Erweiterung nur Logarithmen von Elementen
in $F$ braucht und dass diese Logarithmen in der Stammfunktion nur 
linear vorkommen.
Beim Ableiten müssen die Terme, die Elemente der Erweiterung enthalten,
wieder verschwinden.
Man muss also für exponentielle und algebraische Erweiterungen zeigen,
dass diese beim Ableiten nicht verschwinden können.
Und für logarithmische Erweiterungen muss man zeigen, dass diese genau
dann verschwinden können, wenn sie linear in der Stammfunktion auftreten.

In diesem Abschnitt beweisen wir dies für eine Erweiterung, die von einem
einzigen zusätzlichen Element erzeugt werden.
Dies ist zwar nicht ausreichend für das Integral
\[
\int \frac{dx}{x^2+1}
=
\frac{1}{2i}\log(x+i)
-
\frac{1}{2i}\log(x-i)
=
\arctan x,
\]
welches zwei unabhängige logarithmische Erweiterungen braucht, aber
die Methode, die wir dabei kennenlernen werden, wird in
Abschnitt~\ref{buch:intalg:liouville:subsection:allgemein}
verallgemeinert, so dass sie auch für kompliziertere Integrale
gilt und damit das liouvillesche Prinzip bestätigt.

%
% Ein Ausdruck für die Stammfunktion
%
\subsubsection{Ein Ausdruck für die Stammfunktion}
In diesem Abschnitt ist $F$ ein Funktionenkörper mit Konstantenkörper
$K$ und $f\in F$ eine Funktion.
Es wird angenommen, dass es eine Stammfunktion $g$ von $f$ in einem
Erweiterungskörper $F(\vartheta)$ gibt.
$g\in F(\vartheta)$ bedeutet aber, dass $g$ ein rationaler Ausdruck
in $\vartheta$ ist.
Insbesondere kann $g$ als
\begin{equation}
g
=
a_0(\vartheta)
+
\frac{a(\vartheta)}{b(\vartheta)}
\label{buch:intalg:liouville:spezial:eqn:bruch}
\end{equation}
geschrieben werden, wobei $a_0(\vartheta)$, $a(\vartheta)$ und $b(\vartheta)$
Polynome mit Koeffizienten in $F$ sind.
Ausserdem darf angenommen werden, dass $a(\vartheta)$ und $b(\vartheta)$
teilerfremd sind und dass $\deg a(\vartheta)<\deg b(\vartheta)$ sind,
dass der Bruch in 
\eqref{buch:intalg:liouville:spezial:eqn:bruch}
vollständig gekürzt.

Wir nehmen ausserdem an, dass $F$ und $F(\vartheta)$ den gleichen
Konstantenkörper haben.
Dies kann immer erreicht werden, indem man alle für das Integral
nötigen Konstanten bereits dem Körper $F$ hinzufügt, bevor man beginnt,
um nichtkonstante Funktionen zu erweitern.

Wir müssen zeigen, dass $\vartheta$ nur ein Logarithmus sein kann und dass
$\vartheta$ nur linear in $g$ vorkommen kann.

%
% Der Fall einer logarithmischen Erweiterung
%
\subsubsection{Der Fall einer logarithmischen Erweiterung}

\begin{satz}
Sei $F(\vartheta)$ eine logarithmische Erweiterung eines Funktionenkörpers
$F$.
$F$ und $F(\vartheta)$ haben den gleichen Konstantenkörper.
Weiter sei $f\in F$ und $g \in F(\vartheta)$ eine Stammfunktion.
Dann ist $g=c+d\vartheta$ mit $c,d\in F$.
\end{satz}

\begin{proof}
Der Nenner $b(\vartheta)$ des Ausdrucks 
\eqref{buch:intalg:liouville:spezial:eqn:bruch}
für $g$ kann in die irreduziblen Faktoren
$b_1(\vartheta),\ldots,b_m(\vartheta)$
zerlegt werden.
Weiter dürfen wir annehmen, dass die Faktoren $b_k(\vartheta)$
normierte Polynome sind.
Die zugehörige Partialbruchzerlegung hat daher die Form
\begin{align}
g
&=
a_0(\vartheta)
+
\sum_{k=1}^m
\sum_{i=1}^{r_k}
\frac{a_{ik}(\vartheta)}{b_k(\vartheta)^i},
\notag
\intertext{wobei $a_{ik}(\vartheta)\in F(\vartheta)$ Polynome mit
$\deg a_{ik}(\vartheta) < \deg b_k(\vartheta)$ sind.
Die Ableitung von $g$ ist}
f
=
Dg
&=
Da_0(\vartheta)
+
\sum_{k=1}^m
\sum_{i=1}^{r_k}
D
\frac{a_{ik}(\vartheta)}{b_k(\vartheta)^i}.
\label{buch:intalg:liouville:speziell:eqn:fabgeleitet}
\intertext{Dabei ist wichtig, dass $f\in F$ ist und daher die Erweiterung
$\vartheta$ nicht enthält.
Wir wollen zeigen, dass diese zur Folge hat, dass gar keine Bruchterme
vorkommen können.
Die Ableitung eines einzelnen Terms in der Summe ist}
D
\frac{a_{ik}(\vartheta)}{b_k(\vartheta)^i}
&=
\frac{
Da_{ik}(\vartheta)\cdot b_k(\vartheta)^i
-
a_{ik}(\vartheta)ib_k(\vartheta)^{i-1} Db_k(\vartheta)
}{
b_k(\vartheta)^{2i}
}
\notag
\\
&=
\frac{Da_{ik}(\vartheta)}{b_k(\vartheta)^i}
-
i
\frac{
a_{ik}(\vartheta) Db_k(\vartheta)
}{
b_k(\vartheta)^{i+1}
}.
\label{buch:intalg:liouville:speziell:eqn:abgeleitet}
\end{align}
Die Ableitung $Db_k(\vartheta)$ hat Grad
$\deg Db_k(\vartheta) = \deg b_k(\vartheta)-1$.
Da $b_k(\vartheta)$ irreduzibel ist, sind $Db_k(\vartheta)$
und $b_k(\vartheta)$ teilerfremd.
Die Terme 
\eqref{buch:intalg:liouville:speziell:eqn:abgeleitet}
sind alle vollständig gekürzt.
Die höchste Nennerpotenz von $b_k(\vartheta)$ in
\eqref{buch:intalg:liouville:speziell:eqn:fabgeleitet}
hat den Exponenten $r_k+1$, und sie kommt nur genau einmal
im letzten Summanden vor.
Daher gibt es keinen weiteren Term, mit dem er sich wegheben könnte.
Die Summe in~\eqref{buch:intalg:liouville:speziell:eqn:fabgeleitet}
muss daher verschwinden und $g$ ist nur ein Polynom von $a_0(\vartheta)$.

Da $f=Dg=Da_0(\vartheta)$ nicht von $\vartheta$ abhängen darf, muss
$a_0(\vartheta)$ ein Polynom vom Grad höchstens $1$ sein, also
\[
g
=
a_0(\vartheta) = a_{0,0}+a_{0,1}\vartheta,
\]
wie behauptet.
\end{proof}

%
% Der Fall einer exponentiellen Erweiterung
%
\subsubsection{Der Fall einer exponentiellen Erweiterung}

%
% Der Fall einer algebraischen Erweiterung
%
\subsubsection{Der Fall einer algebraischen Erweiterung}

%
% Beweis im allgemeinen Fall
%
\subsection{Beweis im allgemeinen Fall
\label{buch:intalg:liouville:subsection:allgemein}}



